\documentclass[11pt]{article}
\usepackage[a4paper,margin=26mm]{geometry}
\usepackage{amsmath,amssymb,amsthm}
\usepackage[T1]{fontenc}
\usepackage{xurl}
\usepackage{hyperref}
\hypersetup{colorlinks=true,linkcolor=blue,urlcolor=blue,citecolor=blue,
 pdftitle={The maximal energy of integral circulant graphs of orders p\textasciicircum(2r) 3\textasciicircum(2s+1) and pq\textasciicircum m},
 pdfauthor={Horace Dong}}
\newtheorem{theorem}{Theorem}
\newtheorem{lemma}[theorem]{Lemma}
\newtheorem{corollary}[theorem]{Corollary}
\newtheorem{proposition}[theorem]{Proposition}
\newtheorem{conjecture}[theorem]{Conjecture}
\theoremstyle{remark}
\newtheorem{remark}[theorem]{Remark}
\newcommand{\ICG}{\operatorname{ICG}}
\newcommand{\E}{\mathcal E}
\newcommand{\NA}{\operatorname{NA}}
\title{The maximal energy of integral circulant graphs\\ of orders $p^{2r}3^{2s+1}$ and $pq^m$}
\author{Horace Dong\\ \small Independent researcher\\
\small\texttt{maxwelldhx@gmail.com}\\
\small ORCID: \href{https://orcid.org/0009-0002-7554-7548}{0009-0002-7554-7548}}
\date{September 2026}
\begin{document}
\maketitle
\begin{abstract}
Jiang and Yang determined the maximal energy of integral circulant graphs of
order $n=p^{2r}q^{2s+1}$, where $p,q$ are distinct odd primes with $q\ge5$:
the checkerboard divisor set $\{p^iq^j: i+j \text{ even}\}$ is the unique
maximiser. They noted that their semidefinite method does not extend to
$q=3$. We prove the remaining case: for every prime $p\ge5$ and all $r\ge1$,
$s\ge0$, the checkerboard set is the unique energy maximiser for order
$p^{2r}3^{2s+1}$, with the maximal energy given by the formula of Jiang and
Yang. This settles the maximal-energy problem for opposite-parity exponents
for all pairs of distinct odd primes. The key step is a sharp inequality for
sign matrices under weighted Ramanujan transforms. We prove it with a path
identity and an explicit potential function, and the same argument gives a
semidefinite-free proof of the theorem of Jiang and Yang. For exponents of
equal parity the checkerboard set contains $n$ and is not admissible. We
conjecture the maximal energy in this case, with two maximisers, and prove
the conjecture for the orders $pq^m$ and $p^mq$ with $m$ odd and $p,q$ any
distinct odd primes; for $n=pq$ it also follows from known energy formulas.
With the opposite-parity case, this determines the maximal energy for every
order $pq^m$. Exact computer-assisted certificates for eleven further
exponent pairs show that the conjecture holds for all such orders $p^aq^b$
with $(a+1)(b+1)\le27$. The main theorem, the sign-matrix inequality, the
energy formula for these graphs and the theorem for the orders $pq^m$ and
$p^mq$ are formally verified in Lean~4; the closed-form evaluation of the
maximal energies and the certificates are not formalised. The proofs were found by
AI coding agents working under the author's direction; their role is
described at the end of the paper.
\end{abstract}

\section{Introduction}
The energy of a finite graph is the sum of the absolute values of its
adjacency eigenvalues~\cite{gutman}. For $n\ge2$ and a nonempty set
$\mathcal D$ of proper divisors of $n$, the integral circulant graph
$\ICG_n(\mathcal D)$ has vertex set $\mathbb Z/n\mathbb Z$, two distinct
vertices $u,v$ being adjacent when $\gcd(u-v,n)\in\mathcal D$. These are the
circulant graphs with integral spectrum~\cite{so}. The problem of maximising
the energy of $\ICG_n(\mathcal D)$ over $\mathcal D$ for fixed $n$ has been
solved for prime powers~\cite{sander2013} and studied for other orders; see
the introduction of~\cite{jy} for a survey. Rold\'an~\cite{roldan} conjectured
that for $n=p^2q^3$ the checkerboard set $\{1,p^2,pq,q^2,p^2q^2,pq^3\}$ is the
unique maximiser. Jiang and Yang~\cite{jy} proved much more.

\begin{theorem}[Jiang--Yang {\cite[Theorem~A]{jy}}]\label{thm:jy}
Let $p,q$ be distinct odd primes with $q\ge5$, let $r\ge1$, $s\ge0$ and
$n=p^{2r}q^{2s+1}$. Among all nonempty sets of proper divisors of $n$, the
energy of $\ICG_n(\mathcal D)$ is uniquely maximised by
\[
 \mathcal D^*_{r,s}=\{p^iq^j:\ 0\le i\le 2r,\ 0\le j\le 2s+1,\ i+j\ \text{even}\},
\]
and the maximal energy is $\tfrac12\bigl[n+d_{2r}(p)\,d_{2s+1}(q)\bigr]$, where
$d_k(x)=(2k+1)x^k+4\sum_{j=0}^{k-1}(-1)^{k-j}(j+1)x^j$.
\end{theorem}

They remark that the case $q=3$ is not covered, because their odd-exponent
semidefinite argument fails at $x=3$~\cite[Remark~4.2]{jy}. Our main result
removes this restriction.

\begin{theorem}\label{thm:main}
The conclusion of Theorem~\ref{thm:jy} holds for $q=3$: for every prime
$p\ge5$ and all $r\ge1$, $s\ge0$, the set $\mathcal D^*_{r,s}$ is the unique
energy maximiser among nonempty sets of proper divisors of $n=p^{2r}3^{2s+1}$,
and the maximal energy is $\tfrac12\bigl[n+d_{2r}(p)\,d_{2s+1}(3)\bigr]$.
\end{theorem}

Together with Theorem~\ref{thm:jy} (which allows $p=3$ on the even side),
Theorem~\ref{thm:main} determines the maximal energy and the unique maximiser
for every order $p^{\mathrm{even}}q^{\mathrm{odd}}$ with $p,q$ distinct odd
primes. The case $r=s=1$ was treated earlier. Park~\cite{park} claims a proof of
Rold\'an's conjecture for all distinct odd primes. Schreib~\cite{schreib}
gives a written proof of the case $q=3$ and combines it with~\cite{jy} to
obtain the conjecture.
An arithmetic certificate for $n=27p^2$ was announced in~\cite{tru}. These
works concern only the exponents $(2,3)$, and we are not aware of any earlier
treatment of the case $q=3$ for general $r$ and $s$.

If the two exponents have the same parity, the checkerboard set contains $n$
itself and is not a set of proper divisors, and the maximiser changes. For
this case we state a conjecture (Conjecture~\ref{conj:equal}), with two
maximisers and an explicit maximal energy, and we prove it when one of the
exponents equals~$1$. For $n=p^aq^b$, a set written as $\{p^iq^j:\ \dots\}$
consists of the divisors $p^iq^j$ of $n$ ($0\le i\le a$, $0\le j\le b$) with
the stated property, and $\delta_k(x)=\bigl(x^k(x-1)+2(-1)^k\bigr)/(x+1)$
(Remark~\ref{rem:closed}).

\begin{theorem}\label{thm:onem}
Let $p,q$ be distinct odd primes, let $m\ge1$ be odd, and let $n=pq^m$. The
maximal energy of an integral circulant graph of order $n$ is
\[
 \tfrac12\bigl[n+(3p-4)\,d_m(q)\bigr]-(p-2)\,\delta_m(q),
\]
and it is attained exactly by the sets $\{p^iq^j:\ i+j\ \text{odd}\}$ and
$\{p^iq^j:\ i+j\ \text{even}\}\setminus\{n\}$. Exchanging the names of the
primes, for $n=p^mq$ the maximal energy is
$\tfrac12\bigl[n+(3q-4)\,d_m(p)\bigr]-(q-2)\,\delta_m(p)$, attained exactly by
the two sets with the same description.
\end{theorem}

For $m=1$ the maximal energy is $4(p-1)(q-1)$, attained by $\{p,q\}$ and
$\{1\}$. This case can also be derived from published energy formulas of
Ili\'c~\cite{ilic} and of Ili\'c and Ba\v{s}i\'c~\cite{ilicbasic}
(Remark~\ref{rem:pq}), although we have not found it stated. We are not aware
of an earlier determination of the maximal energy for other orders with
exponents of equal parity. If $m$ is even, the exponents of $pq^m$ have
opposite parity. Theorems~\ref{thm:jy}, \ref{thm:main} and~\ref{thm:onem}
therefore determine the maximal energy, and all maximisers, for every order
$pq^m$ ($m\ge1$) with $p,q$ distinct odd primes. Exact computer-assisted
certificates (Proposition~\ref{prop:cert}) prove the conjecture for every
exponent pair $(a,b)$ with $a,b\ge2$ and $(a+1)(b+1)\le27$; in general it
remains open.

Our proof replaces the semidefinite part of~\cite{jy} by
Theorem~\ref{thm:sign} below, an inequality for sign matrices that holds for
all exponents and parities and all real $p\ge5$, $q\ge3$. By symmetry it also
yields Theorem~\ref{thm:jy} without semidefinite certificates
(Remark~\ref{rem:swap}). Section~\ref{sec:equal} treats exponents of equal
parity, and Section~\ref{sec:onem} proves Theorem~\ref{thm:onem}; the proof
applies the path identity and the potential of Sections~\ref{sec:path}
and~\ref{sec:potential} to the two rows of the sign matrix.

\section{Weighted Ramanujan transforms}
We use the notation of~\cite{jy}. For real $x>1$, an integer $k\ge1$ and
$0\le i,j\le k$, let $T_k(x)=(t_{ij})$ be the symmetric matrix with
\[
t_{ij}=\begin{cases}0,& i+j\le k-2,\\ -x^{k-1}(x-1),& i+j=k-1,\\
x^{2k-i-j-2}(x-1)^2,& i+j\ge k,\end{cases}\qquad (i,j<k),
\]
$t_{ik}=t_{ki}=x^{k-i-1}(x-1)$ for $i<k$, and $t_{kk}=1$. When $x$ is prime,
$t_{ij}=\varphi(x^{k-i})\,c_{x^{k-j}}(x^i)$, where $c_m(u)$ denotes
Ramanujan's sum~\cite[(2.2)--(2.4)]{jy}. The row sums are
$T_k(x)\mathbf 1=x^ke_k$~\cite[(2.5)]{jy}, where $e_k$ is the last standard
basis vector. For a real matrix $Z$ let $\|Z\|_1$ be the sum of the absolute
values of its entries.

\begin{proposition}[{\cite[Proposition~2.2]{jy}}]\label{prop:model}
Let $p\ne q$ be primes, $a,b\ge1$, $n=p^aq^b$, and let $X\in\{0,1\}^{(a+1)
\times(b+1)}$ have $x_{ij}=1$ exactly when $p^iq^j\in\mathcal D$. Then
$\E(\ICG_n(\mathcal D))=\|T_a(p)XT_b(q)^{\mathsf T}\|_1$.
\end{proposition}

Let $s_k=(1,-1,1,\dots,(-1)^k)^{\mathsf T}$ and define
$d_k(x)=\|T_k(x)s_k\|_1$, and put $d_0(x)=1$. By Remark~\ref{rem:closed}
below, $d_k(x)$ agrees with the closed form in Theorem~\ref{thm:jy} for all
real $x>2$; see also~\cite[(3.6)]{jy}.

\section{A path identity}\label{sec:path}
For $w\in\mathbb R^{k+1}$ define $z_{k-1}=w_k$ and, for
$j=k-1,k-2,\dots,0$,
\begin{equation}\label{eq:path}
z_{j-1}=\frac{z_j+(x-1)w_j}{x}.
\end{equation}
Thus $z_j$ is a convex combination of $w_{j+1},\dots,w_k$, and each $z_j$ is a
linear function of $w$.

\begin{lemma}\label{lem:path}
For $w\in\mathbb R^{k+1}$ we have
\[
\|T_k(x)w\|_1=x^k\Bigl(|z_{-1}|+\sum_{j=0}^{k-1}|z_{j-1}-z_j|\Bigr).
\]
\end{lemma}
\begin{proof}
Unwinding~\eqref{eq:path} gives
$z_j=x^{-(k-j-1)}\sum_{i>j}\varphi(x^{k-i})w_i$ with $\varphi(x^0)=1$ and
$\varphi(x^e)=x^{e-1}(x-1)$. The last row of $T_k(x)$ is
$(\varphi(x^{k-i}))_i$, so $(T_k(x)w)_k=x^kz_{-1}$. For $u<k$ put
$j=k-1-u$. Row $u$ has entry $-x^{k-1}(x-1)$ in column $j$, entries
$x^{j}(x-1)\varphi(x^{k-i})$ in columns $i>j$, and zeros elsewhere. Hence
$(T_k(x)w)_u=x^{k-1}(x-1)(z_j-w_j)$, while
$z_{j-1}-z_j=\frac{x-1}{x}(w_j-z_j)$.
\end{proof}

\section{The one-dimensional potential}\label{sec:potential}
For $x>2$ let $\mu_{-1}=1$, $\mu_j=(x-1-\mu_{j-1})/x$ for $j\ge0$, and
$c_j=(x-1)(1+\mu_{j-1})/x$. By induction,
$\mu_j\in\bigl[\frac{x-2}{x},1\bigr]$ for all $j$.

\begin{lemma}\label{lem:potential}
Let $x>2$, let $y\in\{\pm1\}^{k+1}$, and let $z_j$ be defined
by~\eqref{eq:path} with $w=y$. Then $\operatorname{sgn} z_j=y_{j+1}$ and
$\frac{x-2}{x}\le|z_j|\le1$ for $-1\le j\le k-1$, and
\[
\frac{\|T_k(x)y\|_1}{x^k}=\sum_{j=0}^{k-1}c_j+\mu_{k-1}
 -\sum_{j\in\NA(y)}2\mu_j|z_j|,\qquad
\NA(y)=\{0\le j\le k-1:\ y_j=y_{j+1}\}.
\]
Consequently $\|T_k(x)y\|_1\le d_k(x)$, with equality exactly when
$y=\pm s_k$, and $d_k(x)=x^k(\sum_jc_j+\mu_{k-1})$. Moreover
$(T_k(x)s_k)_k>0$.
\end{lemma}
\begin{proof}
We have $z_{k-1}=y_k$, and since $x-1>1\ge|z_j|$, the sign of
$z_{j-1}=(z_j+(x-1)y_j)/x$ is $y_j$ and
$|z_{j-1}|=(x-1\pm|z_j|)/x\in[\frac{x-2}{x},1]$. Put
$f_j=|z_{j-1}-z_j|+\mu_{j-1}|z_{j-1}|-\mu_j|z_j|$. If $y_j=-\operatorname{sgn}z_j$,
then $|z_{j-1}-z_j|=\frac{x-1}{x}(1+|z_j|)$ and
$|z_{j-1}|=\frac{x-1-|z_j|}{x}$; using $x\mu_j=x-1-\mu_{j-1}$ one finds
$f_j=c_j$. If $y_j=\operatorname{sgn}z_j$, the same computation with
$|z_{j-1}-z_j|=\frac{x-1}{x}(1-|z_j|)$ and $|z_{j-1}|=\frac{x-1+|z_j|}{x}$
gives $f_j=c_j-2\mu_j|z_j|$. Summing, the terms $\mu_j|z_j|$ telescope, and
Lemma~\ref{lem:path} gives the identity, since $\mu_{-1}=1$ and
$|z_{k-1}|=1$. Every $j\in\NA(y)$ contributes a positive term
$2\mu_j|z_j|$, and $\NA(y)=\emptyset$ exactly for $y=\pm s_k$. Finally
$(T_k(x)s_k)_k=x^kz_{-1}$ and $\operatorname{sgn}z_{-1}=(s_k)_0=1$.
\end{proof}

\begin{remark}\label{rem:closed}
Put $\delta_k(x)=x^k\mu_{k-1}$. Solving the recursion for $\mu$ gives
$\delta_k(x)=\bigl(x^k(x-1)+2(-1)^k\bigr)/(x+1)$. For $y=s_k$, every step in
the proof of Lemma~\ref{lem:potential} has $y_j=-\operatorname{sgn}z_j$, so
$|z_{j-1}|=(x-1-|z_j|)/x$; starting from $|z_{k-1}|=1=\mu_{-1}$ this gives
$|z_j|=\mu_{k-2-j}$, and in particular
$(T_k(x)s_k)_k=x^kz_{-1}=\delta_k(x)$. Since the potential does not
depend on $k$, Lemma~\ref{lem:potential} gives
$d_{k+1}/x^{k+1}-d_k/x^k=c_k+\mu_k-\mu_{k-1}=2\mu_k$, that is,
$d_{k+1}(x)=x\,d_k(x)+2\delta_{k+1}(x)$, with $d_0=1$. The closed form
$(2k+1)x^k+4\sum_{j=0}^{k-1}(-1)^{k-j}(j+1)x^j$ satisfies the same recursion
and initial value, so it equals $d_k(x)$ for all $k\ge0$.
\end{remark}

\section{A two-variable inequality}
\begin{lemma}\label{lem:cell}
Let $P\ge\mu$ and $0\le\mu\le1$, and put
$h(A,B)=P|A-B|+\mu|B+PA|-(P-\mu)|B|-P(1+\mu)|A|$ for real $A,B$. Then
$h(A,B)\le2\mu\max\{0,|B|-P|A|\}$.
\end{lemma}
\begin{proof}
As $h(-A,-B)=h(A,B)$ we may assume $A\ge0$. If $0\le B\le A$ then
$h=-2(P-\mu)B\le0$. If $B>A$ then $h=2(\mu B-PA)\le2\mu(B-PA)$ because
$\mu\le1$. If $-PA\le B<0$ then $h=0$. If $B<-PA$ then $h=2\mu(|B|-PA)$.
\end{proof}

\begin{lemma}\label{lem:local}
Let $q\ge3$, $b\ge1$, $N=T_b(q)$, $F(w)=\|Nw\|_1$ and $D=d_b(q)$. Let
$P\ge4$ and $\frac35\le\mu\le1$. For all $y\in\{\pm1\}^{b+1}$ and
$\zeta\in[-1,1]^{b+1}$,
\[
\Lambda(y,\zeta):=P\,F(y-\zeta)+\mu\,F(\zeta+Py)-(P-\mu)\,F(\zeta)
\le P(1+\mu)\,D,
\]
with strict inequality unless $y=\pm s_b$.
\end{lemma}
\begin{proof}
Let $\xi_j(w)$ be the path~\eqref{eq:path} with $x=q$, $k=b$, and put
$\tau_{-1}(w)=\xi_{-1}(w)$, $\tau_j(w)=\xi_{j-1}(w)-\xi_j(w)$ for $0\le j<b$.
These are linear, and $F(w)=q^b\sum_j|\tau_j(w)|$ by Lemma~\ref{lem:path}.
With $A_j=\tau_j(y)$ and $B_j=\tau_j(\zeta)$,
\[
\Lambda/q^b=P(1+\mu)F(y)/q^b+\textstyle\sum_jh(A_j,B_j).
\]
By Lemma~\ref{lem:potential} (with $x=q$ and potential $\mu^q_j$), it
suffices to show $\sum_jh(A_j,B_j)\le\sum_{j\in\NA(y)}2P(1+\mu)\mu^q_jx_j$,
where $x_j=|\xi_j(y)|\in[\frac{q-2}{q},1]$, with strict inequality when
$\NA(y)\ne\emptyset$. Since $\xi_j(\zeta)$ is a convex combination of entries
of $\zeta$, we have $|B_{-1}|\le1$ and
$|B_j|=\frac{q-1}{q}|\zeta_j-\xi_j(\zeta)|\le\frac{2(q-1)}{q}$. Also
$|A_{-1}|\ge\frac{q-2}{q}$, and $|A_j|=\frac{q-1}{q}(1+x_j)$ for
$j\notin\NA(y)$, $|A_j|=\frac{q-1}{q}(1-x_j)$ for $j\in\NA(y)$. We use
Lemma~\ref{lem:cell}, noting $P\ge4>\mu$.
\begin{itemize}
\item $j=-1$: $P|A_{-1}|\ge\frac{4(q-2)}{q}\ge\frac43>|B_{-1}|$, so $h\le0$.
\item $j\ge0$, $j\notin\NA(y)$: $P|A_j|\ge\frac{8(q-1)^2}{q^2}\ge
 \frac{2(q-1)}{q}\ge|B_j|$, so $h\le0$.
\item $j\in\NA(y)$: it suffices that $\mu|B_j|<P\beta(x_j)$, where
 $\beta(t)=\mu\frac{q-1}{q}(1-t)+(1+\mu)\mu^q_jt$. As $\beta$ is affine it
 suffices to check the endpoints. For $q=3$, $\beta(1)\ge\frac85\cdot\frac13$
 and $\beta(\frac13)\ge\frac35\cdot\frac49+\frac85\cdot\frac19=\frac49$, so
 $P\beta\ge\frac{16}{9}>\frac43\ge\mu|B_j|$. For general $q\ge3$ the same
 comparison reduces to $P\beta(1)\ge\frac{32}{5}\cdot\frac{q-2}{q}>
 \frac{2(q-1)}{q}$ and $\frac{22}5q^2-\frac{94}5q+\frac{104}5>0$.
\end{itemize}
Summing over $j$ proves the claim.
\end{proof}

\section{The sign-matrix inequality}
\begin{theorem}\label{thm:sign}
Let $p\ge5$ and $q\ge3$ be real, $a,b\ge1$, $M=T_a(p)$ and $N=T_b(q)$. For
every $Y\in\{\pm1\}^{(a+1)\times(b+1)}$,
\[
\|MYN^{\mathsf T}\|_1\le d_a(p)\,d_b(q),
\]
with equality exactly when $Y=\pm s_as_b^{\mathsf T}$.
\end{theorem}
\begin{proof}
Let $\bar y_0,\dots,\bar y_a$ be the rows of $Y$. Define $\zeta_{a-1}=\bar
y_a$ and $\zeta_{k-1}=(\zeta_k+(p-1)\bar y_k)/p\in[-1,1]^{b+1}$. Column $v$ of
$MYN^{\mathsf T}$ is $Mw^{(v)}$ with $w^{(v)}_i=(N\bar y_i)_v$, and by
linearity the path of $w^{(v)}$ is $((N\zeta_k)_v)_k$. Applying
Lemma~\ref{lem:path} with $x=p$ to every column and summing,
\[
\|MYN^{\mathsf T}\|_1=p^a\Bigl(F(\zeta_{-1})+\sum_{k=0}^{a-1}
F(\zeta_{k-1}-\zeta_k)\Bigr),\qquad F(w)=\|Nw\|_1.
\]
Let $\mu_k$ be the potential of Lemma~\ref{lem:potential} for $x=p$ and put
$f_k=F(\zeta_{k-1}-\zeta_k)+\mu_{k-1}F(\zeta_{k-1})-\mu_kF(\zeta_k)$.
Telescoping gives
$\|MYN^{\mathsf T}\|_1/p^a=\sum_kf_k+\mu_{a-1}F(\bar y_a)$. Since
$\zeta_{k-1}-\zeta_k=\frac{p-1}{p}(\bar y_k-\zeta_k)$ and
$p\mu_k=p-1-\mu_{k-1}$, we get $pf_k=\Lambda(\bar y_k,\zeta_k)$ with
$P=p-1\ge4$ and $\mu=\mu_{k-1}\in[\frac{p-2}{p},1]\subseteq[\frac35,1]$.
Lemma~\ref{lem:local} gives $f_k\le c_kd_b(q)$, and
Lemma~\ref{lem:potential} gives $F(\bar y_a)\le d_b(q)$; both are strict
unless the row involved is $\pm s_b$, and $\mu_{a-1}>0$. Hence
$\|MYN^{\mathsf T}\|_1\le p^a(\sum_kc_k+\mu_{a-1})d_b(q)=d_a(p)d_b(q)$.
If equality holds, every row of $Y$ is $\pm s_b$, so $Y=gs_b^{\mathsf T}$ with
$g\in\{\pm1\}^{a+1}$ and $\|MYN^{\mathsf T}\|_1=\|Mg\|_1\,d_b(q)$; by
Lemma~\ref{lem:potential} this equals $d_a(p)d_b(q)$ only for $g=\pm s_a$.
\end{proof}

\begin{remark}\label{rem:swap}
Since $\|MYN^{\mathsf T}\|_1=\|NY^{\mathsf T}M^{\mathsf T}\|_1$, the roles of
the two factors may be exchanged, so Theorem~\ref{thm:sign} holds whenever
$\min(p,q)\ge3$ and $\max(p,q)\ge5$. In particular it contains the sign-matrix
inequality~\cite[Lemma~4.1]{jy} used for Theorem~\ref{thm:jy}, and the
deduction below then proves Theorem~\ref{thm:jy} as well. The semidefinite
argument of~\cite{jy} does not extend to $x=3$~\cite[Remark~4.2]{jy}; the
proof above instead uses the integrality of the signs on both factors.
\end{remark}

\section{Proof of Theorem~\ref{thm:main}}\label{sec:proofmain}
Let $a=2r$, $b=2s+1$, $M=T_a(p)$, $N=T_b(3)$, let $X$ be the divisor matrix
of $\mathcal D$ (so $x_{ab}=0$), and let $Y=2X-J$ with $J$ the all-ones
matrix. By Proposition~\ref{prop:model} and $M\mathbf 1=p^ae_a$,
$N\mathbf1=3^be_b$,
\[
2\E(\ICG_n(\mathcal D))=\|MJN^{\mathsf T}+MYN^{\mathsf T}\|_1
=\|n\,e_ae_b^{\mathsf T}+MYN^{\mathsf T}\|_1\le n+d_a(p)\,d_b(3),
\]
by Theorem~\ref{thm:sign}. For the checkerboard set, $Y^*=s_as_b^{\mathsf T}$,
$x^*_{00}=1$, and $x^*_{ab}=0$ because $a+b$ is odd, so $\mathcal D^*_{r,s}$
is an admissible nonempty set. The $(a,b)$ entry of
$MY^*N^{\mathsf T}=(Ms_a)(Ns_b)^{\mathsf T}$ is positive by
Lemma~\ref{lem:potential}, so it does not cancel against $n$, and
$2\E(\ICG_n(\mathcal D^*_{r,s}))=n+\|Ms_a\|_1\|Ns_b\|_1=n+d_a(p)\,d_b(3)$. If
$\E(\ICG_n(\mathcal D))=\E(\ICG_n(\mathcal D^*_{r,s}))$, then
$\|MYN^{\mathsf T}\|_1=d_a(p)d_b(3)$, so $Y=\pm s_as_b^{\mathsf T}$ by
Theorem~\ref{thm:sign}; since $y_{ab}=-1=(s_as_b^{\mathsf T})_{ab}$, we get
$Y=Y^*$ and $\mathcal D=\mathcal D^*_{r,s}$. \qed

\section{Exponents of equal parity}\label{sec:equal}
Let $p,q$ be distinct odd primes, let $a,b\ge1$ with $a+b$ even, $n=p^aq^b$,
$M=T_a(p)$ and $N=T_b(q)$. For a set $\mathcal D$ of proper divisors of $n$
let $X$ and $Y=2X-J$ be as in Section~\ref{sec:proofmain}, so that
$y_{ab}=-1$, and let $Z=MYN^{\mathsf T}$, with entries $Z_{ij}$. The last rows
of $M$ and $N$ are nonnegative with sums $p^a$ and $q^b$, so $|Z_{ab}|\le n$.
Writing $t^-=\max\{-t,0\}$, the computation of Section~\ref{sec:proofmain}
gives
\begin{equation}\label{eq:G}
G(Y):=2\E(\ICG_n(\mathcal D))-n=\sum_{(i,j)\ne(a,b)}|Z_{ij}|+Z_{ab}
=\|Z\|_1-2(Z_{ab})^-.
\end{equation}

Since $a+b$ is even, the sign matrix $s_as_b^{\mathsf T}$ has entry $+1$ at
$(a,b)$ and violates the proper-divisor condition. The matrix
$-s_as_b^{\mathsf T}$, which corresponds to the set $\{p^iq^j: i+j\text{
odd}\}$, is admissible, and it attains the bound $\|Z\|_1\le d_a(p)\,d_b(q)$
of Theorem~\ref{thm:sign}. However, its entry
$Z_{ab}=-(Ms_a)_a(Ns_b)_b=-\delta_a(p)\,\delta_b(q)$ is negative by
Remark~\ref{rem:closed}, and it cancels against $n$ in the argument of
Section~\ref{sec:proofmain}. So the energy bound
$\tfrac12[n+d_a(p)\,d_b(q)]$ is not attained, and the maximiser changes.
By~\eqref{eq:G}, the set $\{p^iq^j: i+j\text{ odd}\}$ has energy
$\tfrac12[n+d_a(p)\,d_b(q)]-\delta_a(p)\,\delta_b(q)$.

\begin{conjecture}\label{conj:equal}
Let $p,q$ be distinct odd primes, $a,b\ge1$ with $a+b$ even, and
$n=p^aq^b$. The maximal energy of an integral circulant graph of order $n$ is
\[
\tfrac12\bigl[n+d_a(p)\,d_b(q)\bigr]-\delta_a(p)\,\delta_b(q),
\]
attained exactly by the sets $\{p^iq^j: i+j\text{ odd}\}$ and
$\{p^iq^j: i+j\text{ even}\}\setminus\{n\}$.
\end{conjecture}

By the computation above, the first of the two sets always attains the
conjectured value. Since $d_1(x)=3x-4$ and $\delta_1(x)=x-2$,
Theorem~\ref{thm:onem} is the case $\min\{a,b\}=1$ of the conjecture; it is
proved in Section~\ref{sec:onem}. The next result is computer-assisted.

\begin{proposition}\label{prop:cert}
Conjecture~\ref{conj:equal} holds for the exponent pairs $(a,b)=(2,2)$,
$(2,4)$, $(4,2)$, $(2,6)$, $(6,2)$, $(2,8)$, $(8,2)$, $(3,3)$, $(3,5)$, $(5,3)$
and $(4,4)$, and all pairs of distinct odd primes $p,q$.
\end{proposition}

\begin{proof}[Proof (computer-assisted)]
Exchanging the names of the primes transposes $X$ and does not change the
energy, and the list of exponent pairs is closed under transposition. As one
of two distinct odd primes is at least $5$, it therefore suffices to show,
for each listed $(a,b)$ and all real $p\ge5$, $q\ge3$, that every
$Y\in\{\pm1\}^{(a+1)\times(b+1)}$ with $y_{ab}=-1$ satisfies
$G(Y)\le\Theta_{ab}:=d_a(p)\,d_b(q)-2\delta_a(p)\,\delta_b(q)$, with equality
exactly for $Y=-s_as_b^{\mathsf T}$ and $Y=s_as_b^{\mathsf T}-2e_ae_b^{\mathsf
T}$, the sign matrices of the two sets in the conjecture; for real $p,q$,
$G(Y)$ is defined by the middle expression in~\eqref{eq:G}.

Substitute $p=5+s$ and $q=3+t$ with $s,t\ge0$. Then every entry of $Z$ is a
polynomial in $s$ and $t$ with integer coefficients, and so is $\Theta_{ab}$,
because $d_k(x)$ and $\delta_k(x)=x^k+2\sum_{i<k}(-1)^{k-i}x^i$ have integer
coefficients (Remark~\ref{rem:closed}). An entry whose coefficients are all
$\ge0$, or all $\le0$, has a known sign on the whole region. Choose a sign
$\varepsilon_{ij}$ for each of the other entries with $(i,j)\ne(a,b)$, and
let $L$ be the sum of $Z_{ab}$, of the entries of known sign multiplied by
their signs, and of the $\varepsilon_{ij}Z_{ij}$. At every point of the
region, $G(Y)$ is the largest of these linear forms $L$. Hence, if for every
choice of signs $\Theta_{ab}-L$ either vanishes identically or has
nonnegative coefficients and a positive constant term, then
$G(Y)\le\Theta_{ab}$ on the region; equality holds identically if some choice
gives $\Theta_{ab}-L=0$, and $G(Y)<\Theta_{ab}$ everywhere otherwise. For
each listed $(a,b)$, a computer calculation in exact integer arithmetic over
all $2^{(a+1)(b+1)-1}$ admissible sign matrices (up to $2^{26}$) shows that
this condition holds for every $Y$, and that some choice gives
$\Theta_{ab}-L=0$ exactly when $Y$ is one of the two matrices above. The
programs are described in Section~\ref{sec:verif}.
\end{proof}

Together with Theorem~\ref{thm:onem}, Conjecture~\ref{conj:equal} therefore
holds whenever $\min\{a,b\}=1$ or $(a+1)(b+1)\le27$. It is open for all other
exponent pairs of equal parity; the smallest are $(3,7)$, $(2,10)$ and their
transposes.

\begin{remark}\label{rem:pq}
For $n=pq$, Theorem~\ref{thm:onem} gives the maximal energy $4(p-1)(q-1)$,
attained exactly by $\{p,q\}$ and $\{1\}$. This can also be read off from
known formulas. The graph $\ICG_{pq}(\{1\})$ is the unitary Cayley graph,
whose energy $4\varphi(pq)$ was computed by Ili\'c~\cite{ilic}, and
$\E(\ICG_{pq}(\{p,q\}))=4\varphi(pq)$ by~\cite[Theorem~4.2]{ilicbasic}. The
graphs $\ICG_{pq}(\{p\})$ and $\ICG_{pq}(\{1,p\})$ are $p$ disjoint copies of
$K_q$ and the complete $q$-partite graph with parts of size~$p$
(see also~\cite[Theorem~4.1]{ilicbasic}); both have energy $2p(q-1)$, and
symmetrically for $\{q\}$ and $\{1,q\}$. Finally $\ICG_{pq}(\{1,p,q\})$ is the
complete graph $K_{pq}$, of energy $2(pq-1)$. We have not found the maximal
energy for $n=pq$ stated explicitly in the literature. Le and
Sander~\cite[Theorem~2.1 and Corollary~2.1]{lesander} determined the maximum
$\tilde E_{\max}(n)$ of the energy over \emph{multiplicative} divisor sets and
showed $E_{\max}(n)\ge\tilde E_{\max}(n)$. For $n=pq$ equality holds, but for
$n=pq^3$ it fails: by direct computation, $E_{\max}(5\cdot3^3)=584>512=\tilde
E_{\max}(5\cdot3^3)$ and $E_{\max}(3\cdot5^3)=1632>1472=\tilde
E_{\max}(3\cdot5^3)$, and in both cases the maximisers are the
non-multiplicative sets of Theorem~\ref{thm:onem}.
\end{remark}

\begin{remark}\label{rem:near}
Other sets come close to the maximum. Let $(a,b)=(2,2)$, $q=3$ and
$\mathcal D=\{1,3,p^2,3p^2,9p\}$, so that $Y=s_2y^{\mathsf T}$ with
$y=(1,1,-1)^{\mathsf T}$ and $Z=(Ms_2)(Ny)^{\mathsf T}$. The path of $y$ is
$z_1=-1$, $z_0=\frac13$, $z_{-1}=\frac79$, and $\NA(y)=\{0\}$, so
Lemma~\ref{lem:potential} gives $\|Ny\|_1=d_2(3)-2=23$, and
$Z_{22}=9\,\delta_2(p)\,z_{-1}>0$. By~\eqref{eq:G}, $G(Y)=23\,d_2(p)$, while
the maximum of $G$ is $25\,d_2(p)-10\,\delta_2(p)=G(Y)+4(p-3)$
(Proposition~\ref{prop:cert}). So the energy of $\mathcal D$ is smaller than
the maximal energy $62p^2-90p+40$ only by $2(p-3)$.
\end{remark}

\section{Proof of Theorem~\ref{thm:onem}}\label{sec:onem}
Let $p,q,m$ and $n=pq^m$ be as in Theorem~\ref{thm:onem}; the case $n=p^mq$
follows by exchanging the names of the primes. Put $x=q$, $P=x-1$ and
$R=p-1$. As $p$ and $q$ are distinct odd primes, either $x\ge5$ and $R\ge2$,
or $x=3$ and $R\ge4$; only these inequalities are used. In both cases
\begin{equation}\label{eq:RP}
R+1<P(R-1).
\end{equation}
Let $\mu_j$ and $c_j$ be the potential of Lemma~\ref{lem:potential} for $x$,
and put $\rho=\mu_{m-1}$, $\hat c_j=c_j$ for $0\le j<m$, $\hat c_m=\rho$ and
$\hat D=\sum_{j=0}^m\hat c_j$. By Lemma~\ref{lem:potential} and
Remark~\ref{rem:closed}, $\hat D=d_m(x)/x^m$ and $\rho=\delta_m(x)/x^m$. For
$w\in\mathbb R^{m+1}$ let $z_j(w)$, $-1\le j\le m-1$, be the
path~\eqref{eq:path} with $k=m$, and let $\nu(w)=\|T_m(x)w\|_1/x^m$. For a
sign vector $y$ put $\sigma(y)=\sum_{j\in\NA(y)}\mu_j|z_j(y)|$, so that
$\nu(y)=\hat D-2\sigma(y)$ by Lemma~\ref{lem:potential}. For a statement $S$,
$[S]$ is $1$ if $S$ holds and $0$ otherwise.

For $0\le\mu\le1$ let $h_\mu$ be the function $h$ of Lemma~\ref{lem:cell}
with $P=x-1$. Besides Lemma~\ref{lem:cell} we use three values, which follow
from the definition and from the case analysis in the proof of
Lemma~\ref{lem:cell}: $h_\mu(0,B)=2\mu|B|$; $h_\mu(A,B)=-2(P-\mu)|B|$ if
$AB>0$ and $|B|\le|A|$; and $h_\mu(A,B)=0$ if $AB\le0$ and $|B|\le P|A|$.

Let $\mu_\infty=P/(x+1)$. As $x\mu_\infty=P-\mu_\infty$, the recursion for
$\mu$ gives $\mu_j-\mu_\infty=(-1/x)^{j+1}(1-\mu_\infty)$, where
$1-\mu_\infty>0$. Hence, for $j\ge0$,
\begin{equation}\label{eq:muinf}
\mu_j>\mu_\infty\ \text{if $j$ is odd},\qquad
\mu_j<\mu_\infty\ \text{if $j$ is even},\qquad
\mu_0\le\mu_j\le\mu_1.
\end{equation}
In particular $\rho<\mu_\infty$, because $m-1$ is even.

\begin{lemma}\label{lem:realpot}
For every $w\in\mathbb R^{m+1}$,
\[
\nu(w)=\sum_{j=0}^m\hat c_j|w_j|
+\frac1x\sum_{j=0}^{m-1}h_{\mu_{j-1}}\bigl(w_j,z_j(w)\bigr).
\]
\end{lemma}
\begin{proof}
Write $z_j=z_j(w)$. By~\eqref{eq:path}, $z_{j-1}-z_j=\frac Px(w_j-z_j)$ and
$xz_{j-1}=z_j+Pw_j$. Together with $x\mu_j=P-\mu_{j-1}$ this gives, for
$0\le j<m$,
\[
|z_{j-1}-z_j|+\mu_{j-1}|z_{j-1}|-\mu_j|z_j|
=\tfrac1x\bigl[h_{\mu_{j-1}}(w_j,z_j)+P(1+\mu_{j-1})|w_j|\bigr].
\]
Summing over $j$, the potential terms telescope to
$\mu_{-1}|z_{-1}|-\mu_{m-1}|z_{m-1}|=|z_{-1}|-\rho|w_m|$, and
Lemma~\ref{lem:path} gives the claim.
\end{proof}

\begin{lemma}\label{lem:noncorner}
For $y\in\{\pm1\}^{m+1}$ we have $\nu(y)-|z_{-1}(y)|\le\hat D-\rho$, with
equality exactly when $y=\pm s_m$.
\end{lemma}
\begin{proof}
By Lemma~\ref{lem:path}, the left side is $\sum_{j<m}|z_{j-1}-z_j|$, where
$z_j=z_j(y)$. Let $\mu'_{-1}=0$ and $\mu'_j=(P-\mu'_{j-1})/x$ for $j\ge0$, so
that $\mu'_j\in(0,P/x]$ for $j\ge0$. The computation in the proof of
Lemma~\ref{lem:potential} uses only this recursion of the potential, so it
gives
\[
|z_{j-1}-z_j|+\mu'_{j-1}|z_{j-1}|-\mu'_j|z_j|
=c'_j-2\mu'_j|z_j|\,[j\in\NA(y)],\qquad c'_j=P(1+\mu'_{j-1})/x.
\]
Summing over $j<m$ and using $\mu'_{-1}=0$ and $|z_{m-1}|=1$,
\[
\sum_{j<m}|z_{j-1}-z_j|=\sum_{j<m}c'_j+\mu'_{m-1}
-2\sum_{j\in\NA(y)}\mu'_j|z_j|.
\]
All terms of the last sum are positive, so the left side is maximal exactly
when $\NA(y)=\emptyset$, that is, when $y=\pm s_m$. For $y=s_m$ it equals
$\nu(s_m)-|z_{-1}(s_m)|=\hat D-\rho$, by Lemma~\ref{lem:potential} and
Remark~\ref{rem:closed}.
\end{proof}

\begin{lemma}\label{lem:paritydefect}
Let $y\in\{\pm1\}^{m+1}$ with $y_0=y_m$. Then $(3R-1)\,\sigma(y)>(R-1)\rho$.
\end{lemma}
\begin{proof}
As $m$ is odd and $y_0=y_m$, $y\ne\pm s_m$, so $\NA(y)$ is nonempty; let
$J=\max\NA(y)$. Since $y$ alternates on $J+1,\dots,m$, the recursion
$|z_{j-1}|=(P-|z_j|)/x$ from the proof of Lemma~\ref{lem:potential} holds for
$J<j\le m-1$, and starting from $|z_{m-1}|=1=\mu_{-1}$ it gives
$|z_J(y)|=\mu_{m-2-J}$. Hence $\sigma(y)\ge\mu_J\mu_{m-2-J}$. If $J=m-1$,
this is $\rho$, and $(3R-1)\rho>(R-1)\rho$. Otherwise $J$ and $m-2-J$ are
nonnegative and have opposite parity, since their sum is odd, so
$\sigma(y)>\mu_0\mu_\infty$ by~\eqref{eq:muinf}. As
$(3R-1)(x-2)-(R-1)x=2R(x-3)+2>0$, we have $(3R-1)\mu_0>R-1$, and therefore
$(3R-1)\sigma(y)>(3R-1)\mu_0\mu_\infty\ge(R-1)\mu_\infty\ge(R-1)\rho$.
\end{proof}

\begin{proof}[Proof of Theorem~\ref{thm:onem}]
Let $\mathcal D$ be a set of proper divisors of $n$, and let $u$ and $v$ be
the two rows of its sign matrix $Y$, so that $v_m=-1$. Since
$T_1(p)=\bigl(\begin{smallmatrix}-R&R\\R&1\end{smallmatrix}\bigr)$, the rows
of $Z=T_1(p)YT_m(x)^{\mathsf T}$ are $-R\,T_m(x)w^1$ and $T_m(x)w^2$, where
$w^1=u-v$ and $w^2=Ru+v$. By the proof of Lemma~\ref{lem:path},
$Z_{1m}=x^mz_{-1}(w^2)$, so~\eqref{eq:G} gives $G(Y)=x^m\Phi$ with
\[
\Phi=R\,\nu(w^1)+\nu(w^2)-2\bigl(z_{-1}(w^2)\bigr)^-.
\]
As $d_1(p)=3R-1$ and $\delta_1(p)=R-1$, the claimed maximal energy is
$\frac12(n+x^m\Theta)$ with $\Theta=(3R-1)\hat D-2(R-1)\rho$. We show that
$\Phi\le\Theta$, with equality exactly for $Y=s_1s_m^{\mathsf
T}-2e_1e_m^{\mathsf T}$ and $Y=-s_1s_m^{\mathsf T}$. These are the sign
matrices of the two sets in the theorem, which are nonempty sets of proper
divisors of $n$ because $1+m$ is even; so the theorem follows.

Let $\mathcal A=\{j:\ u_j=v_j\}$. For $j\in\mathcal A$ we have $w^1_j=0$ and
$|w^2_j|=R+1$; for $j\notin\mathcal A$ we have $|w^1_j|=2$ and
$|w^2_j|=R-1$. So $R|w^1_j|+|w^2_j|$ is $R+1$ on $\mathcal A$ and $3R-1$ off
$\mathcal A$, and Lemma~\ref{lem:realpot}, applied to $w^1$ and $w^2$, gives
\begin{equation}\label{eq:Phi}
\Theta-\Phi=2(R-1)\Bigl(\sum_{j\in\mathcal A}\hat c_j-\rho\Bigr)-H
+2\bigl(z_{-1}(w^2)\bigr)^-,
\end{equation}
where $H$ is the sum of the $2m$ \emph{cells}
$\frac Rx\,h_{\mu_{j-1}}\bigl(w^1_j,z_j(w^1)\bigr)$ and
$\frac1x\,h_{\mu_{j-1}}\bigl(w^2_j,z_j(w^2)\bigr)$, $0\le j<m$. Recall that
$z_j(w)$ is a convex combination of $w_{j+1},\dots,w_m$.
\begin{itemize}
\item[(a)] For the $w^2$-cells, $|z_j(w^2)|\le R+1<P(R-1)\le P|w^2_j|$
 by~\eqref{eq:RP}, so they are $\le0$ by Lemma~\ref{lem:cell}.
\item[(b)] For the $w^1$-cells with $j\notin\mathcal A$,
 $|z_j(w^1)|\le2\le P|w^1_j|$, so they are $\le0$.
\item[(c)] A $w^1$-cell with $j\in\mathcal A$ equals
 $\frac Rx h_{\mu_{j-1}}(0,z_j(w^1))=\frac{2R}x\mu_{j-1}|z_j(w^1)|$.
\end{itemize}
For $j\in\mathcal A$ with $j<m$ let
\[
s_j=2(R-1)c_j-\tfrac{2R}x\,\mu_{j-1}|z_j(w^1)|
\ \ge\ \tfrac2x\bigl[(R-1)P(1+\mu)-2R\mu\bigr],\qquad \mu=\mu_{j-1},
\]
where we used $|z_j(w^1)|\le2$. The bracket is affine in $\mu$, and it equals
$(R-1)P>0$ at $\mu=0$ and $2[R(P-1)-P]>0$ at $\mu=1$; hence $s_j>0$. Let
$S=\sum_{j\in\mathcal A,\,j<m}s_j$, let $C\ge0$ be minus the sum of the cells
in (a) and~(b), and let $K=2(z_{-1}(w^2))^-\ge0$. Then~\eqref{eq:Phi} reads
\begin{equation}\label{eq:budget}
\Theta-\Phi=S+C+K-2(R-1)\rho\,[m\notin\mathcal A].
\end{equation}
Moreover $|z_0(w^2)|\le R+1<P|w^2_0|$, so $z_{-1}(w^2)=(z_0(w^2)+Pw^2_0)/x$
has the sign of $w^2_0=Ru_0+v_0$, which is the sign of $u_0$. Thus $K=0$ if
and only if $u_0=1$.

\emph{Case 1: $m\in\mathcal A$.} Then $\Theta-\Phi=S+C+K\ge0$. Suppose that
equality holds. Then $S=0$, so $\mathcal A=\{m\}$; $K=0$, so $u_0=1$; and
$C=0$. Now $w^1_j=2u_j$ for $j<m$ and $w^1_m=0$. For $j\le m-2$ we have
$|Pw^1_{j+1}|=2P>2\ge|z_{j+1}(w^1)|$, so
$z_j(w^1)=(z_{j+1}(w^1)+Pw^1_{j+1})/x$ has the sign of $u_{j+1}$, and
$|z_j(w^1)|\le2=|w^1_j|$. If $u_j=u_{j+1}$ for some $j\le m-2$, the
$w^1$-cell at $j$ is therefore $-\frac{2R}x(P-\mu_{j-1})|z_j(w^1)|<0$,
contradicting $C=0$. So $u$ alternates on $0,\dots,m-1$ and $u_0=1$; together
with $u_m=v_m=-1$ and $v_j=-u_j$ for $j<m$, this means
$Y=s_1s_m^{\mathsf T}-2e_1e_m^{\mathsf T}$, as $m$ is odd. Conversely, for
this $Y$ we have $\mathcal A=\{m\}$, $u=s_m$ and $S=K=0$. The $w^1$-cells
vanish: for $j\le m-2$, $w^1_j=2u_j$ and $z_j(w^1)$ have opposite signs and
$|z_j(w^1)|\le P|w^1_j|$, and $z_{m-1}(w^1)=w^1_m=0$. The $w^2$-cells vanish
too: $w^2_i$ has the sign of $u_i$ for every $i$; by~\eqref{eq:RP},
$z_j(w^2)=(z_{j+1}(w^2)+Pw^2_{j+1})/x$ has the sign of $w^2_{j+1}$ for
$j\le m-2$, and $z_{m-1}(w^2)=w^2_m$; since $u$ alternates, $w^2_j$ and
$z_j(w^2)$ have opposite signs, and $|z_j(w^2)|\le P|w^2_j|$ by~(a). Hence
$\Phi=\Theta$.

\emph{Case 2: $m\notin\mathcal A\ne\emptyset$.} Then $u_m=1$ and
$\mathcal A\subseteq\{0,\dots,m-1\}$. By~\eqref{eq:budget} it suffices to
find $k\in\mathcal A$ with $s_k>2(R-1)\rho$. Since
$\rho<\mu_\infty=P/(P+2)$, the lower bound for $s_k$ shows that this holds
if
\begin{equation}\label{eq:case2}
(R-1)P+\mu(P+2)\bigl[R(P-2)-P\bigr]\ge0,\qquad \mu=\mu_{k-1}.
\end{equation}
If $x\ge5$, then $P\ge4$ and $R(P-2)-P\ge P-4\ge0$, so~\eqref{eq:case2}
holds for every $k\in\mathcal A$. If $x=3$, then~\eqref{eq:case2} reads
$2(R-1)\ge8\mu$. If some $k\in\mathcal A$ has $k\ge1$, then
$\mu_{k-1}\le\mu_1=\frac59$ by~\eqref{eq:muinf}, and
$2(R-1)\ge6>\frac{40}9$. It remains to treat $x=3$ and $\mathcal A=\{0\}$.
Here $\mu_\infty=\frac12$ and $c_0=\frac43$, and since $z_0(w)$ depends only
on $w_1,\dots,w_m$, we have $z_0(w^1)=2z_0(u)$ and
$s_0=\frac23\bigl[4(R-1)-2R|z_0(u)|\bigr]$. If $m=1$, then $|z_0(u)|=1$,
$\rho=\mu_0=\frac13$ and $s_0=\frac23(2R-4)>\frac23(R-1)=2(R-1)\rho$. Let
$m\ge3$. If $|z_0(u)|\le\frac{5(R-1)}{4R}$, then $s_0\ge R-1>2(R-1)\rho$.
Otherwise $|z_0(u)|>\frac{15}{16}$. If $u_1\ne u_2$, then
$u_1=-\operatorname{sgn}z_1(u)$ by Lemma~\ref{lem:potential}, and
$|z_0(u)|=(2-|z_1(u)|)/3\le\frac59$; hence $u_1=u_2$. Then
$1\notin\mathcal A$, and the arguments $w^1_1=2u_1$ and $z_1(w^1)=2z_1(u)$ of
the $w^1$-cell at $j=1$ have the same sign, with $|z_1(w^1)|\le|w^1_1|$. This
cell equals
$-\frac{2R}3(P-\mu_0)|z_1(w^1)|=-\frac{20R}9|z_1(u)|\le-\frac{20R}{27}$, so
$C\ge\frac{20R}{27}$. With $s_0\ge\frac23(2R-4)$ and $2(R-1)\rho<R-1$,
\[
\Theta-\Phi\ge\tfrac23(2R-4)+\tfrac{20R}{27}-2(R-1)\rho
>\tfrac{56R-72}{27}-(R-1)=\tfrac{29R-45}{27}>0.
\]

\emph{Case 3: $\mathcal A=\emptyset$}, that is, $v=-u$. Then $u_m=1$,
$w^1=2u$, $w^2=(R-1)u$ and $\Phi=(3R-1)\nu(u)-2(R-1)(z_{-1}(u))^-$. If
$u_0=1$, then $z_{-1}(u)>0$ by Lemma~\ref{lem:potential}, and
$\Phi=(3R-1)(\hat D-2\sigma(u))<\Theta$ by Lemma~\ref{lem:paritydefect},
since $u_0=u_m$. If $u_0=-1$, then $(z_{-1}(u))^-=|z_{-1}(u)|$ and
\[
\Phi=(R+1)\,\nu(u)+2(R-1)\bigl(\nu(u)-|z_{-1}(u)|\bigr)
\le(R+1)\hat D+2(R-1)(\hat D-\rho)=\Theta
\]
by Lemmas~\ref{lem:potential} and~\ref{lem:noncorner}, with equality exactly
when $u=\pm s_m$, that is, $u=-s_m$. This is $Y=-s_1s_m^{\mathsf T}$.

The three cases cover all $Y$, which proves the theorem. For $m=1$ the value
is $4(p-1)(q-1)$, since $d_1(q)=3q-4$ and $\delta_1(q)=q-2$.
\end{proof}

\section{Verification}\label{sec:verif}
The proofs above are complete, except that the proof of
Proposition~\ref{prop:cert} rests on the computer calculation described in the
paragraph on certificates below. The following computations are independent
cross-checks.
\begin{itemize}
\item An exhaustive search over all divisor sets, in exact integer
 arithmetic, confirms Theorem~\ref{thm:main}, including the value of the
 maximum, in the following cases:
 \begin{itemize}
 \item exponent pairs $(4,1)$, $(6,1)$, $(8,1)$, $(2,5)$, $(4,3)$, for the
  primes $5\le p\le31$;
 \item $(10,1)$, $(12,1)$, $(2,7)$, $(6,3)$, for $p\in\{5,7,11,13,17\}$;
 \item $(14,1)$, $(4,5)$, $(2,9)$, for $p\in\{5,7,11\}$ (the largest case has
  $2^{29}$ divisor sets).
 \end{itemize}
\item Theorem~\ref{thm:sign} was checked in exact arithmetic over all sign
 matrices of every shape with $(a+1)(b+1)\le25$, for $(p,q)=(5,3)$, $(7,3)$
 and $(5,4)$.
\item Each lemma was tested separately in exact arithmetic: by adversarial
 searches for small $b$, and by structured random tests of
 Lemma~\ref{lem:local} up to $b=20$.
\item Proposition~\ref{prop:model} was compared, in floating-point
 arithmetic, with spectra of the actual circulant matrices. The fast Fourier
 transform was used for all divisor sets of the orders $n=75$, $147$, $363$,
 $675$, $1323$, $1875$ and $6075$. Dense diagonalisation was used for all
 divisor sets when $n\le675$ and for samples of them when $n=1323$, $1875$ and
 $6075$.
\item Theorem~\ref{thm:onem} was confirmed by exhaustive searches over all
 divisor sets: in exact arithmetic, with the eigenvalues computed from
 Ramanujan sums and independently of Proposition~\ref{prop:model}, for $35$
 orders $pq^m$ with $m\in\{1,3,5,7\}$ up to $n=234375$; and with
 floating-point spectra from the fast Fourier transform for ten orders up to
 $n=9375$.
\item The inequality $\Phi\le\Theta$ of Section~\ref{sec:onem}, with its
 equality cases, was checked in exact arithmetic over all pairs of sign
 vectors for $m\le11$ at several parameter pairs $(x,R)$. Each intermediate
 claim of its proof was checked separately for $m\le7$, including
 non-integer parameters.
\item An exhaustive search over all divisor sets confirms
 Conjecture~\ref{conj:equal} for ten ordered pairs of odd primes and each of
 the exponent pairs $(1,1)$, $(2,2)$, $(1,3)$, $(3,1)$, $(3,3)$, $(2,4)$,
 $(4,2)$, $(4,4)$, $(1,5)$ and $(3,5)$. These cases are now covered by
 Theorem~\ref{thm:onem} and Proposition~\ref{prop:cert}.
\end{itemize}

\paragraph{Certificates for Proposition~\ref{prop:cert}.}
The certificate is implemented in Python, in exact integer arithmetic, in two
versions that share the routines building the polynomial matrices. A
vectorised version uses $64$-bit integers, with an asserted bound that
excludes overflow; it first applies the cheaper sufficient test in which the
absolute value of each entry of undetermined sign is bounded by the
polynomial with the absolute values of its coefficients, and branches over
signs only when this test fails. A pure-Python version branches over the
signs of the undetermined entries for every sign matrix. The vectorised
version certified all eleven exponent pairs; the pure-Python version gave the
same result for $(2,2)$, $(2,4)$, $(4,2)$ and $(3,3)$. An independent
re-implementation in SymPy, which builds $T_k(x)$ from the Ramanujan-sum
formula and recomputes the target as $G(-s_as_b^{\mathsf T})$, checking it
against the closed form, confirmed $(2,2)$, $(2,4)$, $(4,2)$ and $(3,3)$.
The seven larger exponent pairs were run only with the vectorised version.
The polynomial matrices were compared with the exact matrices $T_k(x)$,
$k\le8$, at sample points. As negative controls, the certificate is rejected
when the target is lowered by~$1$, and, for $(a,b)=(1,1)$ and $(2,2)$, when
the region is enlarged to $p,q\ge3$: at $p=q=3$ further sign matrices attain
the target. Proposition~\ref{prop:cert} is not formalised.

\paragraph{Formal verification.}
Theorem~\ref{thm:sign}, and Theorems~\ref{thm:main} and~\ref{thm:onem} with
$d_k(x)$ taken as $\|T_k(x)s_k\|_1$ and $\delta_k(x)$ as $(T_k(x)s_k)_k$, are
verified in Lean~4 (version
4.33.1, Mathlib commit \texttt{0df444a3}). The development starts from the
definition of the graph and of its energy:
\begin{itemize}
\item the adjacency matrix of $\ICG_n(\mathcal D)$ has entry~$1$ exactly when
 $\gcd((j-i)\bmod n,n)\in\mathcal D$;
\item the energy is the sum of the absolute values of its eigenvalues, as
 given by Mathlib's spectral theorem.
\end{itemize}
It proves:
\begin{itemize}
\item Lemmas~\ref{lem:path}, \ref{lem:potential}, \ref{lem:cell}
 and~\ref{lem:local}, and Theorem~\ref{thm:sign}, for
 real $p\ge5$, $q\ge3$ and all $a,b\ge1$, together with the exchanged form of
 Remark~\ref{rem:swap};
\item Proposition~\ref{prop:model} for all distinct primes $p,q$ and all
 exponents, from the eigenvalues of circulant matrices;
\item the conclusion of Theorem~\ref{thm:main} and the value of the maximal
 energy, with $d_k(x)$ defined as $\|T_k(x)s_k\|_1$;
\item the conclusion of Theorem~\ref{thm:jy} for all pairs of distinct odd
 primes, which includes Theorem~\ref{thm:main} and the case $q\ge5$;
\item Theorem~\ref{thm:onem}, as the theorem
 \texttt{ICGEqualParity.corollaryA}. For distinct primes $p,q\ge3$ and odd
 $m$, with $n=pq^m$, $\mathcal D^-=\{p^iq^j: i+j\text{ odd}\}$ and
 $\mathcal D^+=\{p^iq^j: i+j\text{ even}\}$, it states that $\mathcal D^-$ and
 $\mathcal D^+\setminus\{n\}$ are sets of proper divisors of $n$; that
 $2\E(\ICG_n(\mathcal D^-))=n+(3p-4)\|T_m(q)s_m\|_1-2(p-2)(T_m(q)s_m)_m$;
 that $\E(\ICG_n(\mathcal D^+\setminus\{n\}))=\E(\ICG_n(\mathcal D^-))$; and
 that every set $\mathcal D$ of proper divisors of $n$ satisfies
 $\E(\ICG_n(\mathcal D))\le\E(\ICG_n(\mathcal D^-))$, with equality only if
 $\mathcal D=\mathcal D^-$ or $\mathcal D=\mathcal D^+\setminus\{n\}$. The case
 $n=p^mq$ is the same statement with the primes exchanged. The inequality
 $\Phi\le\Theta$ of Section~\ref{sec:onem}, strict except for the two
 extremal sign matrices, is proved for all real $x,R$ with $x\ge5$, $R\ge2$
 or $x=3$, $R\ge4$.
\end{itemize}
The closed forms of $d_k(x)$ and $\delta_k(x)$ in Remark~\ref{rem:closed} are
proved above but not formalised; in the Lean statements, $d_m(q)$ and
$\delta_m(q)$ appear as $\|T_m(q)s_m\|_1$ and $(T_m(q)s_m)_m$.
The whole development was rebuilt from a clean directory; the files for
Theorem~\ref{thm:onem} were rebuilt in the same way, together with all files
they import. The recorded axiom
audits of the main results (the command \texttt{\#print axioms}) report only
\texttt{propext}, \texttt{Classical.choice} and \texttt{Quot.sound}; no native
evaluation is used. The code and logs are available at~\cite{repo}.

\section*{Tool and computational resource disclosure}
This work was carried out between 25 and 26 September 2026 in a workflow
directed by the author, using the AI coding agents Anthropic Claude Code
(model Claude Opus~5.5) and OpenAI Codex (desktop application; models
\texttt{gpt-6-astra} and \texttt{gpt-6-sol}). The author chose the research
programme and directed the agents.
\begin{itemize}
\item A Claude Code agent found the proof of Theorem~\ref{thm:sign} and of
 Theorem~\ref{thm:main}, formulated Conjecture~\ref{conj:equal}, and wrote
 the verification programs.
\item Another Claude Code agent wrote the Lean formalisation.
\item The proof was checked line by line by a second Claude Code agent and by
 Codex, and the agents searched the literature.
\item The results on exponents of equal parity, Theorem~\ref{thm:onem} and
 Proposition~\ref{prop:cert}, were found by a Claude Code research agent,
 and Theorem~\ref{thm:onem} was checked by an independent Claude referee
 agent.
\end{itemize} Claude Code drafted this text. The programs and logs are
available at~\cite{repo}.
The author is not a professional mathematician. The author has read the
paper and, using a side-by-side table, compared the statements of the Lean
theorems with the theorems stated here. The correctness of
Theorems~\ref{thm:main}, \ref{thm:sign} and~\ref{thm:onem} rests on the
written proofs and on the Lean formalisation described in
Section~\ref{sec:verif}; Proposition~\ref{prop:cert} rests on the computer
calculation described there; the other independent checks are additional
cross-checks. In addition, the arguments for Theorems~\ref{thm:main}
and~\ref{thm:sign} were checked by AI agents of two different model families,
the argument for Theorem~\ref{thm:onem} was checked by an independent Claude
agent, and the formal statements were reviewed by an agent that did not write
them. AI systems are not authors of this paper; the author takes
full responsibility for its content.

\begin{thebibliography}{99}
\bibitem{repo}
H.~Dong, \emph{ai4math-results}, repository,
\url{https://github.com/Horace-Maxwell/ai4math-results}; archived at Zenodo,
\url{https://doi.org/10.5281/zenodo.22970254}.
\bibitem{gutman}
I.~Gutman, The energy of a graph, \emph{Ber.\ Math.-Statist.\ Sekt.\
Forsch.\ Graz} \textbf{103} (1978), 1--22.
\bibitem{ilic}
A.~Ili\'c, The energy of unitary Cayley graphs, \emph{Linear Algebra Appl.}
\textbf{431} (2009), 1881--1889.
\bibitem{ilicbasic}
A.~Ili\'c and M.~Ba\v{s}i\'c, New results on the energy of integral circulant
graphs, \emph{Appl.\ Math.\ Comput.} \textbf{218} (2011), 3470--3482.
\bibitem{lesander}
T.~A. Le and J.~W. Sander, Extremal energies of integral circulant graphs via
multiplicativity, \emph{Linear Algebra Appl.} \textbf{437} (2012), 1408--1421,
\url{https://doi.org/10.1016/j.laa.2012.04.012}.
\bibitem{jy}
J.~Jiang and C.~Yang, Energy maximisation for integral circulant graphs with
opposite-parity exponents, arXiv:2608.29523v1, 30 August 2026.
\bibitem{park}
S.~Park, \emph{Exact energy maximisation for integral circulant graphs of
order $p^2q^3$}, manuscript posted 21 August 2026,
\url{https://github.com/coshaman/papers/tree/2f730393e3d5bed4d9641d7d79b1bdc57fbff05d}.
\bibitem{roldan}
D.~G. Rold\'an, Graph energy maximisation for integral circulant graphs of
order $n=p^2q^3$, arXiv:2604.09491v1, 10 April 2026.
\bibitem{sander2013}
J.~W. Sander and T.~Sander, The exact maximal energy of integral circulant
graphs with prime power order, \emph{Contrib.\ Discrete Math.} \textbf{8}
(2013), 19--40.
\bibitem{schreib}
J.~A. Schreib, \emph{The Roldan universality conjecture}, repository, 5
September 2026,
\url{https://github.com/jamesschreib/roldan-universality-conjecture/tree/95bb002688c65948a15cc9380e8e1f08f8bd6801}.
\bibitem{so}
W.~So, Integral circulant graphs, \emph{Discrete Math.} \textbf{306} (2006),
153--158.
\bibitem{tru}
The Omega Institute, \emph{Rold\'an integral-circulant energy: exceptional
$q=3$ case}, issue \#8338 of the \texttt{trureturing} repository, opened 16
September 2026,
\url{https://github.com/the-omega-institute/trureturing/issues/8338}.
\end{thebibliography}
\end{document}
